Seja \(\alpha\in\mathbb R\backslash \{ 0 \}\) e suponha que \(F\) e \(G\) são aplicações lineares (operadores) de \(\mathbb R^n\) em \(\mathbb R^n\) satisfazendo \(F\circ G - G\circ F=\alpha F\).
\begin{itemize}

\item[(a)] Mostre que para todo \(k\in\mathbb N\) tem-se \(F^k\circ G-G\circ F^k=\alpha kF^k\).

\item[(b)] Mostre que existe \(k\geq 1\) tal que \(F^k=0\).
\end{itemize}

